% Part: first-order-logic
% Chapter: introduction
% Section: syntax

\documentclass[../../../include/open-logic-section]{subfiles}

\begin{document}

\olfileid{fol}{int}{syn}

% SEGMENT OLP-0141-S01
\olsection{தொடரமைப்பு}

எந்தக் குறித்தொடர்கள் முதல்தர தருக்கத்தின் !!{sentence}s ஆகக் கணக்கிடப்படும்
என்பதை முதலில் துல்லியமாக்க வேண்டும். இதைப் பின்னர் செய்வோம்; இப்போது
எடுத்துக்காட்டுகள் வழியாக மட்டும் தொடர்வோம். முதல்தர தருக்கத்தின் அடிப்படைக்
கட்டுமானக் கூறுகளான சொற்கோவை இரண்டு பகுதிகளாகப் பிரிகிறது. குறிப்பிட்டவற்றைப்
பற்றிக் கூறவும் குறிப்பிட்டவற்றைச் சுட்டவும் நாம் பயன்படுத்தும் குறிகள் முதல்
பகுதி. !!{constant}s ஐப் பயன்படுத்திப் பொருள்களைச் சுட்டுகிறோம்; நாம் சுட்டும்
பொருள்களைப் பற்றி !!{predicate}s ஐப் பயன்படுத்திக் கூறுகிறோம். எடுத்துக்காட்டாக,
ஒரு தனிப் பொருளைச் சுட்ட~$\Obj a$ ஐ !!a{constant} ஆகப் பயன்படுத்தி, பின்னர்
!!{sentence}~$\Atom{\Obj P}{\Obj a}$ ஐப் பயன்படுத்தி அதைப் பற்றிய ஒன்றைக்
கூறலாம். ``$\Obj a$'' மற்றும் ``$\Obj P$'' என்பவற்றுக்கான பொருள்களை மனத்தில்
கொண்டிருந்தால், $\Atom{\Obj P}{\Obj a}$ ஐ ஓர் இயல்மொழி வாக்கியமாகப்
படிக்கலாம் (முறையிய தருக்கத்தை முதன்முதலில் கற்றபோது அவ்வாறே செய்திருக்கவும்
வாய்ப்புள்ளது). முதல்தர தருக்கத்தின் இத்தகைய எளிய !!{sentence}s கிடைத்ததும்,
சொற்கோவையின் இரண்டாம் பகுதியான தருக்கக் குறிகளைப் (இணைப்பிகள் மற்றும்
அளவையடைகள்) பயன்படுத்தி இன்னும் சிக்கலானவற்றைக் கட்டமைக்கலாம். எடுத்துக்காட்டாக,
$(\Atom{\Obj P}{\Obj a} \land \Atom{\Obj Q}{\Obj b})$ அல்லது
$\lexists[\Obj x][\Atom{\Obj P}{\Obj x}]$ போன்ற கோவைகளை உருவாக்கலாம்.

% SEGMENT OLP-0141-S02
கடுமையான மேல்தருக்கவியல் ஆய்வுக்குத் தேவையான பொருண்மையியல் வரையறைகளையும்
நமது !!{derivation} முறைமைகளின் விதிகளையும் வழங்க, முதல்தர தருக்கத்தின்
!!a{sentence} ஆக எது கணக்கிடப்படும் என்பதற்குத் துல்லியமான வரையறையை முதலில்
தர வேண்டும். அடிப்படைக் கருத்தைப் புரிந்துகொள்வது எளிது: !!{predicate}s மற்றும்
!!{constant}s ஐ மட்டும் கொண்டு~$\Atom{\Obj P}{\Obj a}$ போன்ற சில எளிய
!!{sentence}s ஐ உருவாக்கலாம். பின்னர் இணைப்பிகளையும் அளவையடைகளையும்
பயன்படுத்தி அவற்றிலிருந்து இன்னும் சிக்கலானவற்றை உருவாக்கலாம். ஆனால் அவ்வாறு
சிக்கலான !!{sentence}s ஐ உருவாக்க அனுமதிக்கும் விதிகள் துல்லியமாக யாவை?
இவற்றைக் குறிப்பிடாவிட்டால், ``முதல்தர தருக்கத்தின் !!{sentence}'' என்பதைப்
போதுமான துல்லியத்துடன் வரையறுக்கவில்லை. இதில் சில சிக்கல்கள் உள்ளன. சரியான
சரியான !!{sentence}s ஆகக் கணக்கிடப்படும் குறித்தொடர்களைத் தேர்ந்தெடுப்பது
முதலாவது. \emph{எல்லா}
!!{sentence}s ஐயும் பற்றிய கணித நிறுவல்களை வழங்கக்கூடியவாறு இதைச் செய்வது
இரண்டாவது. இறுதியாக, ``$!A$ இலுள்ள ஒவ்வொரு~$\Obj x$ ஐயும்~$\Obj b$ ஆல்
மாற்றுக'' என்பது போன்ற, !!{sentence}s மீதான சில அடிப்படைச் செயல்களுக்கும்
துல்லியமான வரையறைகள் தர வேண்டும். முதல்தர தருக்கத்திலுள்ள அளவையடைகளும்
!!{variable}s உம் இதை எவ்வாறு செய்ய வேண்டும் என்பதை முற்றிலும் வெளிப்படையாக
விடாததே சிக்கல். எடுத்துக்காட்டாக, $\lexists[\Obj x][\Atom{\Obj P}{\Obj a}]$
!!a{sentence} ஆகக் கணக்கிடப்பட வேண்டுமா? $\lexists[\Obj x][\lexists[\Obj
x][\Atom{\Obj P}{\Obj x}]]$ பற்றி என்ன? ``$(\Atom{\Obj P}{\Obj x} \land
\lexists[\Obj x][\Atom{\Obj P}{\Obj x}])$ இல்~$\Obj x$ ஐ~$\Obj b$ ஆல்
மாற்றுக'' என்பதன் விளைவு என்னவாக இருக்க வேண்டும்?

\end{document}
